\section{Sharp adaptation on full interval boxes}
\label{app:budget-adaptation}

This appendix proves Proposition~\ref{prop:histogram-adaptation} for full
interval boxes with separate countable partitions and arbitrary bounded
Borel centers. Throughout, $N=2n$, $q=n^{-1/2}$, $A=a+s^2$, $B=b+t^2$,
and all model and sequential-M1 conventions are those of
Appendix~\ref{app:model}. The covariate law is arbitrary and unknown.
The estimator is fixed for each $(\Gamma,\rho,n)$ before choosing any
of the four certificates or a compatible law, and attains $q+AB$.

\subsection{Public centers and separate represented groupings}

For each $\nu\in\{\mu,\pi\}$, let $(C_{\nu j})_{j\ge1}$ be a
countable Borel partition, and let $c_\nu$ be a bounded Borel function.
The two partitions can differ, and a center need not be constant on
its partition cells. Consider the full interval boxes
\begin{equation}
 \begin{aligned}
 \cG_\nu&=
 \left\{c_\nu+\sum_j\theta_j\ind_{C_{\nu j}}:
       \theta_j\in[-h_{\nu j},h_{\nu j}]\text{ for every }j\right\},\\
 h_{\nu j}&\ge0,\qquad |c_\nu(x)|+h_{\nu j}\le B_{\mathcal G}
               \quad(x\in C_{\nu j}).
 \end{aligned}
 \label{eq:adaptation-boxes}
\end{equation}
Each class contains \emph{every independent} coordinate choice displayed in
\eqref{eq:adaptation-boxes}; inclusion in a box would not suffice.
These classes admit sequential-M1 cores: approximate the normalized
coordinates on the first $k$ cells by rational numbers in $[-1,1]$, and
use zero normalized coordinates elsewhere. Any supplied sequential-M1
core may be used in the construction.

Neither the centers nor the partitions are required as additional inputs.
The original cores recover both centers and widths by
\begin{equation}
 \begin{aligned}
 c_\nu(x)&=\tfrac12\left\{\sup_{\ell\ge1}g_{\nu,\ell}(x)
                         +\inf_{\ell\ge1}g_{\nu,\ell}(x)\right\},\\
 h_\nu(x)&=\tfrac12\left\{\sup_{\ell\ge1}g_{\nu,\ell}(x)
                         -\inf_{\ell\ge1}g_{\nu,\ell}(x)\right\}.
 \end{aligned}
 \label{eq:adaptation-midpoint}
\end{equation}
These countable extrema are Borel and finite. Pointwise sequential M1
transfers the full-class extrema $c_\nu(x)\pm h_{\nu j}$ to the core,
which proves the equality. Define a separate equivalence relation for
each side using its \emph{centered} core:
\begin{equation}
 i\sim_\nu k\quad\Longleftrightarrow\quad
 g_{\nu,\ell}(X_i)-c_\nu(X_i)
   =g_{\nu,\ell}(X_k)-c_\nu(X_k)
 \quad\text{for every }\ell\ge1.
 \label{eq:adaptation-equivalence}
\end{equation}
Each relation is Borel. Equality of all centered core evaluations
implies equality of every centered full-class evaluation, by sequential
M1. Two different cells on side $\nu$ are separated whenever either
has positive width: vary that coordinate while holding the other fixed.
Thus its groups are its occupied partition cells, except that its
zero-width cells merge. Every centered public function on that side
vanishes on the merged cells. No common refinement is taken.

Set the public constants
\begin{equation}
 \alpha=\min\left\{\frac18,\frac{v_-}{8B_0^2}\right\}>0,
 \qquad m=\lfloor\alpha N\rfloor.
 \label{eq:adaptation-trimming-fraction}
\end{equation}
For each side, let $J_\nu$ contain the $m$ singleton rows with largest
widths $h_\nu(X_i)$, or every singleton row if fewer than $m$ exist.
Resolve equal widths by increasing sample index. Let $\mathsf A_\nu$
have block $I_k-\mathbf1_k\mathbf1_k'/k$ on a group of size $k\ge2$,
block $[0]$ on a row in $J_\nu$, and block $[1]$ on every other
singleton. Then
\begin{equation}
 \mathsf A_\nu'=\mathsf A_\nu,\qquad \mathsf A_\nu^2=\mathsf A_\nu,
 \qquad 0\preceq\mathsf A_\nu\preceq I_N,\qquad
 (\mathsf A_\nu)_{ii}\ge\tfrac12\ind_{\{i\notin J_\nu\}}.
 \label{eq:adaptation-projection}
\end{equation}
Every matrix entry is Borel: if
$m_{\nu i}=\sum_k\ind_{\{i\sim_\nu k\}}$, it is
$(\mathsf A_\nu)_{ik}=\ind_{\{i=k\}}\ind_{\{i\notin J_\nu\}}-
 \ind_{\{i\sim_\nu k\}}\ind_{\{m_{\nu i}\ge2\}}/m_{\nu i}$.
Also, $(\mathsf A_\nu)_{ik}\le0$ when $i\ne k$, and
$\|\mathsf A_\nu z\|_\infty\le2\|z\|_\infty$.
For a function $f$, write $f_X=(f(X_1),\ldots,f(X_N))'$,
$\|z\|_N^2=z'z/N$, and $\|z\|_{N,1}=N^{-1}\sum_i|z_i|$.
For a public function $g_\nu=c_\nu+z_\nu$, its centered part obeys
$|z_\nu|\le h_{\nu j}$ on $C_{\nu j}$, and
$\mathsf A_\nu(z_\nu)_X$ vanishes
on repeated groups and trimmed rows, retaining its original values on
the remaining singleton rows. In particular $|J_\nu|\le\alpha N$.

\subsection{Localized complexity controls singleton amplitudes}

The occupancy bound below controls the total singleton amplitude.
Trimming converts this control into a squared residual bound of order
$r^4+N^{-1}$.

\begin{lemma}[An occupancy consequence of the full-box certificate]
\label{lem:adaptation-occupancy}
Fix either class in \eqref{eq:adaptation-boxes}, suppress its index
(including in $\mathsf A=\mathsf A_\nu$), and
let $r\in[0,1]$ be its stochastic upper certificate in (A.1). Set
$p_j=P_X(C_j)$ and
\[
 K_\rho=\max\left\{\frac{3\sqrt3}{2}c_{\rm rad},
                         27B_{\mathcal G}c_{\rm rad}^2\right\}.
\]
Then, for every $n\ge1$, including $r=0$,
\begin{equation}
 \sum_jp_jh_j(1-p_j)^{2n-1}
 \le \sum_jp_jh_j e^{-np_j}\le K_\rho r^2.
 \label{eq:adaptation-singleton-mass}
\end{equation}
Consequently every fixed centered public function $z$ satisfies
\begin{equation}
 \E\|\mathsf A z_X\|_{N,1}\le K_\rho r^2,
 \qquad
 \E\|\mathsf A z_X\|_N^2\le B_{\mathcal G}K_\rho r^2.
 \label{eq:adaptation-residual-moments}
\end{equation}
\end{lemma}

\begin{proof}
Put $\lambda_j=np_j$ and $L=\sum_jp_jh_j e^{-\lambda_j}$. Independent
interval coordinates imply that $\Delta\cG$ contains the damped sign cube
\[
 f_\sigma=2\sum_j\sigma_jh_j e^{-\lambda_j/2}\ind_{C_j},
 \qquad \sigma_j\in\{-1,1\}.
\]
Each such function has squared population norm at most
$4B_{\mathcal G}L$. For the sample of size $n$ used in the complexity
certificate, let $S_j=\sum_{i=1}^n\epsilon_i\ind_{\{X_i\in C_j\}}$.
Then
\[
 \E S_j^2=\lambda_j,\qquad
 \E S_j^4=np_j+3n(n-1)p_j^2\le\lambda_j+3\lambda_j^2.
\]
H\"older's inequality in the form
$\E S_j^2\le(\E|S_j|)^{2/3}(\E S_j^4)^{1/3}$ gives, for
$\lambda_j>0$,
\[
 \E|S_j|\ge\frac{(\E S_j^2)^{3/2}}{(\E S_j^4)^{1/2}}
 \ge\frac{\lambda_j}{\sqrt{1+3\lambda_j}};
\]
the lower bound is also valid when $\lambda_j=0$. Maximizing the signs
coordinatewise produces the measurable random variable
$2n^{-1}\sum_jh_je^{-\lambda_j/2}|S_j|$, which is bounded above by
the full localized supremum at radius $2\sqrt{B_{\mathcal G}L}$.
Since $e^{\lambda/2}/\sqrt{1+3\lambda}\ge1/\sqrt3$ for
$\lambda\ge0$, monotonicity of outer expectation and Tonelli's theorem
give a complexity lower bound $2L/\sqrt3$.

For $r>0$, evaluate (A.1) at
$R=r\vee2\sqrt{B_{\mathcal G}L}$ to obtain
\[
 \frac{2L}{\sqrt3}\le3c_{\rm rad}r
       \max\{r,2\sqrt{B_{\mathcal G}L}\}.
\]
If the maximum is $r$, this bounds $L$ by
$(3\sqrt3/2)c_{\rm rad}r^2$; otherwise, squaring after division by
$\sqrt L$ gives $L\le27B_{\mathcal G}c_{\rm rad}^2r^2$.
If $r=0$ and $L>0$, the same cube at its positive containing radius
contradicts the zero-budget requirement in (A.1), so $L=0$.
No radius-zero quotient is used.

Since $2n-1\ge n$, $(1-p_j)^{2n-1}\le e^{-np_j}$. The expected
empirical amplitude on cells occupied exactly once in the full sample
is exactly
\[
 \E\left[\frac1N\sum_{i=1}^Nh(X_i)
          \ind_{\{\#\{k:X_k\in C(X_i)\}=1\}}\right]
 =\sum_jp_jh_j(1-p_j)^{2n-1},
\]
where $h(x)=h_j$ on $C_j$ and $C(x)$ denotes the cell of $x$. Merged zero-width cells
contribute zero. This proves the first residual bound; the second
follows pointwise from $|z|^2\le B_{\mathcal G}|z|$ on the surviving
singleton groups. The unknown probabilities are used only in this
proof.
\end{proof}

\begin{lemma}[Squared residual bound after trimming]
\label{lem:adaptation-trimmed-residual}
Under Lemma~\ref{lem:adaptation-occupancy}, put
\[
 H=\frac1N\sum_{i=1}^Nh(X_i)
       \ind_{\{\text{the own cell of }X_i\text{ occurs once}\}},
 \qquad
 R_N(r)=\alpha^{-1}
       \sqrt{K_\rho^2r^4+\frac{2B_{\mathcal G}^2}{N}}.
\]
Then
\begin{equation}
 \E H^2\le K_\rho^2r^4+\frac{2B_{\mathcal G}^2}{N},
 \qquad
 \E\|\mathsf A z_X\|_N^2\le R_N(r)^2
 \label{eq:adaptation-trimmed-moments}
\end{equation}
for every fixed centered public function $z$.
\end{lemma}

\begin{proof}
If $N_j$ denotes the number of observations in cell $j$, then
$H=N^{-1}\sum_jh_j\ind_{\{N_j=1\}}$. Replacing one observation changes
only its source and destination cells' contributions, each by at most
its width. Thus the replacement changes $H$ by at most
$2B_{\mathcal G}/N$. Efron--Stein and
Lemma~\ref{lem:adaptation-occupancy} give
\[
 \operatorname{Var}(H)\le2B_{\mathcal G}^2/N,
 \qquad \E H\le K_\rho r^2,
\]
proving the first assertion. Zero-width cells contribute zero, so merging
them in the represented equivalence relation does not change $H$.

If any singleton remains after trimming, let $h_{\max}$ be its largest
width. The $m$ trimmed rows and a remaining row attaining $h_{\max}$ all
have widths at least $h_{\max}$. Therefore
\[
 (m+1)h_{\max}\le NH,\qquad
 h_{\max}\le\frac{NH}{m+1}\le\frac H\alpha,
\]
because $m+1>\alpha N$. This also covers $m=0$. If every singleton was
trimmed, the residual is zero. Otherwise, $\mathsf A z_X$ vanishes on
all repeated groups and trimmed rows, and has absolute entries at most
$h_{\max}$ on its remaining rows. Hence
$\|\mathsf A z_X\|_N^2\le H^2/\alpha^2$, proving the second assertion.
At $r=0$, $\E H=0$ also implies $H=0$ almost surely.
\end{proof}

\subsection{One stabilized rule for the two partitions}

Put $\mathsf B=\mathsf A_\pi\mathsf A_\mu$. This matrix need not be
symmetric or positive semidefinite. Its operator norm is at most one,
and the off-diagonal entries of each factor are nonpositive. Thus
\begin{equation}
 \mathsf B_{ii}
 =\sum_k(\mathsf A_\pi)_{ik}(\mathsf A_\mu)_{ki}
 \ge\frac14\ind_{\{i\notin J_\pi\cup J_\mu\}},
 \qquad
 \mathsf B_{ii}=0\quad(i\in J_\pi\cup J_\mu).
 \label{eq:adaptation-product-matrix}
\end{equation}
For sample vectors $T,Y,c_\mu,c_\pi$, define
\begin{align}
 D&=(T-c_\pi)'\mathsf B T/N,&
 U&=(T-c_\pi)'\mathsf B(Y-c_\mu)/N,\nonumber\\
 S_\mu&=T'\mathsf A_\mu T/N,&
 V_\mu&=T'\mathsf A_\mu(Y-c_\mu)/N.
 \label{eq:adaptation-affine-moments}
\end{align}

\begin{proposition}[All-four-certificate adaptation on full boxes]
\label{prop:affine-histogram-adaptation}
For every fixed sequential-M1 pair of the form
\eqref{eq:adaptation-boxes}, the Borel rule
\begin{equation}
 \widehat\beta_{\rm box}=
 \begin{cases}
 \Pi_{[-3,3]}(U/D),&D\ge v_-/32,\\
 \Pi_{[-3,3]}\{V_\mu/(S_\mu\vee(v_-/4))\},&D<v_-/32
 \end{cases}
 \label{eq:adaptation-affine-rule}
\end{equation}
uses only $(\Gamma,\rho,n)$ and satisfies, with a constant depending
only on $\rho$, simultaneously for every $(a,s,b,t)\in[0,1]^4$ and
every $P\in\mathcal P_{\rho,n,a,s,b,t}(\Gamma)$,
\begin{equation}
 \E_P|\widehat\beta_{\rm box}-\beta(P)|
 \le C_\rho\{q+(a+s^2)(b+t^2)\}.
 \label{eq:adaptation-sharp-risk}
\end{equation}
This includes every $n\ge1$, every zero-budget boundary, and every
regularity tuple with a nonempty compatible-law class.
\end{proposition}

\begin{proof}
Write $\mu_X=(\mu_0)_X$, $\pi_X=(\pi_0)_X$, and define the proof
quantities
\[
 d_\mu=\mu_X-c_\mu,\qquad d_\pi=\pi_X-c_\pi,\qquad
 F_\nu=\|\mathsf A_\nu d_\nu\|_N,\qquad J=\E(F_\mu F_\pi).
\]
For any $\eta>0$, take fixed approximation comparators
$g_\mu=c_\mu+z_\mu$, $g_\pi=c_\pi+z_\pi$ with errors at most
$a+\eta$, $b+\eta$. Contraction and Minkowski's inequality, applied in
$L_2(P^N)$ to their empirical residual norms, give
\[
 (\E F_\mu^2)^{1/2}\le a+\eta+R_N(s),
 \qquad
 (\E F_\pi^2)^{1/2}\le b+\eta+R_N(t).
\]
The rule does not use these comparators. Letting $\eta\downarrow0$
handles unattained approximation infima and yields
\begin{equation}
 J\le\{a+R_N(s)\}\{b+R_N(t)\}.
 \label{eq:adaptation-affine-moments-bound}
\end{equation}

\paragraph{Primary bias and noise.}
The conditional-mean part of the primary score satisfies
\begin{equation}
 |d_\pi'\mathsf B d_\mu|/N
 =|\langle\mathsf A_\pi d_\pi,\mathsf A_\mu d_\mu\rangle_N|
 \le F_\pi F_\mu.
 \label{eq:adaptation-bias}
\end{equation}
Its full decomposition is
\[
 U-\beta_0D=d_\pi'\mathsf B d_\mu/N
       +u'\mathsf B d_\mu/N
       +(T-c_\pi)'\mathsf B\varepsilon/N.
\]
Conditional on $X_{1:N}$, independence and centering of the $u_i$ give
\[
 \E\{(u'\mathsf B d_\mu/N)^2\mid X_{1:N}\}
 \le B_0^2\|\mathsf B d_\mu\|^2/N^2
 \le B_0^2(B_0+B_{\mathcal G})^2/N.
\]
Conditional on $(X_{1:N},T_{1:N})$, the independent centered
$\varepsilon_i$ give the same bound for the other noise term, using
$\|\mathsf B'(T-c_\pi)\|\le\|T-c_\pi\|$. Consequently
\begin{equation}
 \E|U-\beta_0D|
 \le J+2B_0(B_0+B_{\mathcal G})/\sqrt N.
 \label{eq:adaptation-score}
\end{equation}
On the primary branch the clipped loss is at most
$32|U-\beta_0D|/v_-$.

\paragraph{Denominators under average conditional variance.}
Put $v_i=\E(u_i^2\mid X_i)$,
$\overline v=N^{-1}\sum_i v_i$, and
$\Omega=\operatorname{diag}(v_i)$. Since $0\le v_i\le B_0^2$,
$\E\overline v\ge v_-$, and $|J_\nu|\le\alpha N$,
\begin{align}
 \Pp(\overline v<3v_-/4)&\le\frac{16B_0^4}{Nv_-^2},
 \label{eq:adaptation-average-variance}\\
 \operatorname{tr}(\mathsf B\Omega)/N
 &\ge\tfrac14(\overline v-2\alpha B_0^2),\nonumber\\
 \operatorname{tr}(\mathsf A_\mu\Omega)/N
 &\ge\tfrac12(\overline v-\alpha B_0^2).
 \label{eq:adaptation-trimmed-traces}
\end{align}
In particular, on $\overline v\ge3v_-/4$,
\begin{equation}
 \operatorname{tr}(\mathsf B\Omega)/N\ge v_-/8,\qquad
 \E(S_\mu\mid X_{1:N})\ge5v_-/16.
 \label{eq:adaptation-denominator-means}
\end{equation}
The second inequality also uses
$\pi_X'\mathsf A_\mu\pi_X\ge0$. Neither calculation assumes a
pointwise lower bound on $v_i$.

For any symmetric covariate-measurable matrix $H$, independence and
centering give
\begin{align*}
 \operatorname{Var}(u'Hu\mid X_{1:N})
 &=\sum_iH_{ii}^2\operatorname{Var}(u_i^2\mid X_i)
       +4\sum_{i<k}H_{ik}^2v_iv_k\\
 &\le2B_0^4\|H\|_{\rm F}^2.
\end{align*}
The random part of $NS_\mu$ is
$u'\mathsf A_\mu u+2u'\mathsf A_\mu\pi_X$. The two variances are at
most $2B_0^4N$ and $4B_0^4N$. Applying
$\operatorname{Var}(Q+L)\le2\operatorname{Var}Q+
2\operatorname{Var}L$, which does not assume vanishing third moments,
therefore gives
\begin{equation}
 \operatorname{Var}(S_\mu\mid X_{1:N})\le12B_0^4/N.
 \label{eq:adaptation-denominator}
\end{equation}

For the nonsymmetric cross moment, put
\[
 d_X=\E(D\mid X_{1:N})
 =\operatorname{tr}(\mathsf B\Omega)/N
       +d_\pi'\mathsf B\pi_X/N.
\]
Let $H=(\mathsf B+\mathsf B')/2$ and
$\ell_X=\mathsf B\pi_X+\mathsf B'd_\pi$. The random part of $ND$ is
$u'Hu+u'\ell_X$, with
\[
 \|H\|_{\rm op}\le1,\qquad \|H\|_{\rm F}^2\le N,\qquad
 \|\ell_X\|\le(2B_0+B_{\mathcal G})\sqrt N.
\]
The same calculation yields
\begin{equation}
 \operatorname{Var}(D\mid X_{1:N})\le C_D/N,\qquad
 C_D=4B_0^4+2B_0^2(2B_0+B_{\mathcal G})^2.
 \label{eq:adaptation-cross-denominator}
\end{equation}
Define
\[
 \mathcal E=
 \{\overline v<3v_-/4\}\cup
 \{S_\mu<v_-/4\}\cup
 \{|D-d_X|>v_-/32\}.
\]
On $\overline v\ge3v_-/4$, the conditional-mean gap in the second
event is at least $v_-/16$. Conditional Chebyshev and the preceding
bounds, counting the low-$\overline v$ event only once, give
\begin{equation}
 \Pp(\mathcal E)\le
 \frac{3088B_0^4+1024C_D}{Nv_-^2}.
 \label{eq:adaptation-exceptional-cost}
\end{equation}
Clipping limits its contribution to loss to six times this probability.

\paragraph{Fallback bias is controlled by the same product.}
On $\mathcal E^c$, choosing fallback implies $d_X<v_-/16$.
Equation~\eqref{eq:adaptation-denominator-means} then forces
\[
 v_-/16\le|d_\pi'\mathsf B\pi_X|/N
 \le B_0F_\pi.
\]
Thus the remaining fallback event is contained in
$\{F_\pi\ge\kappa\}$, where $\kappa=v_-/(16B_0)$.

The fallback score decomposes as
\[
 V_\mu-\beta_0S_\mu
 =\pi_X'\mathsf A_\mu d_\mu/N
       +u'\mathsf A_\mu d_\mu/N
       +T'\mathsf A_\mu\varepsilon/N.
\]
Its bias is pointwise at most $B_0F_\mu$, and conditioning as above
bounds its total expected absolute noise by
$B_0(2B_0+B_{\mathcal G})/\sqrt N$. Therefore
\[
 \E\!\left[B_0F_\mu
       \ind_{\{\text{fallback}\}\cap\mathcal E^c}\right]
 \le B_0\E\{F_\mu\ind_{\{F_\pi\ge\kappa\}}\}
 \le(B_0/\kappa)J.
\]
Absolute noise terms can be bounded without their event indicators;
branch selection need not be independent of the scores. Divide by
$S_\mu\ge v_-/4$ on $\mathcal E^c$, and combine with the primary
branch and exceptional-event bounds to obtain
\begin{align}
 \E|\widehat\beta_{\rm box}-\beta_0|
 \le{}&
 \left(\frac{32}{v_-}+\frac{64B_0^2}{v_-^2}\right)J
 +\frac{6(3088B_0^4+1024C_D)}{Nv_-^2}\nonumber\\
 &+\left\{\frac{64B_0(B_0+B_{\mathcal G})}{v_-}
       +\frac{4B_0(2B_0+B_{\mathcal G})}{v_-}\right\}N^{-1/2}.
 \label{eq:adaptation-explicit-risk}
\end{align}
Let
$C_*=\max\{1,K_\rho/\alpha,\sqrt2B_{\mathcal G}/\alpha\}$.
Then $R_N(r)\le C_*(r^2+N^{-1/2})$, and
\eqref{eq:adaptation-affine-moments-bound}, with all four budgets at
most one, gives
\[
 J\le C_*^2\{(a+s^2)(b+t^2)+5N^{-1/2}\}.
\]
Since $N=2n$, this proves \eqref{eq:adaptation-sharp-risk}.

Countable extrema, countable centered-core equalities, finite group
sizes, width comparisons with the specified tie rule, matrix operations,
and clipping are Borel. The trimming fraction uses only $\rho$; no
certificate, comparator, or covariate probability enters the rule.
The floor in \eqref{eq:adaptation-trimming-fraction} and
Lemma~\ref{lem:adaptation-trimmed-residual} include all small samples
and zero budgets. This proves simultaneous adaptation.
\end{proof}

With explicitly evaluable own-partition labels, centers, and widths,
the same construction uses finite block operations and sorting.
The original sequential-M1 interface establishes Borel existence,
without asserting finite-query computation for an arbitrary
representation. Full independent coordinate choices are hypotheses of
the occupancy argument; the proposition does not cover arbitrary
represented public classes.

Block residualization is classical \citep{cattaneo2018alternative}.
The present guarantee uses the full-box certificate to choose which
singleton rows to retain and controls the resulting cross-moment fallback;
it does not claim a new partial-regression principle.

\subsection{Matching lower bound over full-box pairs}

Let $\GammaSet^{\rm box}_{\rho,n}(a,s,b,t)$ be the members of
$\GammaSet_{\rho,n}(a,s,b,t)$ whose class-pair projection has the
full-box form \eqref{eq:adaptation-boxes}. Define the restricted
outer-pair minimax risk by
\begin{equation}
 \mathcal R^{\rm box}_{\rho,n}(a,s,b,t)
 =\sup_{\Gamma\in\GammaSet^{\rm box}_{\rho,n}(a,s,b,t)}
   \inf_{\delta_\Gamma\ {\rm Borel}}
   \sup_{P\in\mathcal P_{\rho,n,a,s,b,t}(\Gamma)}
       \E_P|\delta_\Gamma(Z_{1:2n})-\beta(P)|.
 \label{eq:adaptation-box-minimax-risk}
\end{equation}
The infimum here can use the declared budgets. The upper rule in
Proposition~\ref{prop:affine-histogram-adaptation} has the stronger
property of being fixed before all four budgets.

\begin{theorem}[Restricted full-box minimax rate]
\label{thm:full-box-minimax}
On the strict lower-endpoint slack stratum $0<v_-<B_0^2$, there are
constants $c_{{\rm box},\rho},C_{{\rm box},\rho}>0$ such that
\begin{equation}
 c_{{\rm box},\rho}\{q+(a+s^2)(b+t^2)\}
 \le \mathcal R^{\rm box}_{\rho,n}(a,s,b,t)
 \le C_{{\rm box},\rho}\{q+(a+s^2)(b+t^2)\}
 \label{eq:adaptation-box-minimax-rate}
\end{equation}
for every $n\ge1$ and every $(a,s,b,t)\in[0,1]^4$. For each fixed
$(\rho,n,s,t)$, the lower bound is witnessed by one represented full-box
pair and one covariate marginal that work for every $(a,b)\in[0,1]^2$.
It is not a lower bound for every fixed full-box pair. At
$v_-=B_0^2$, the restricted minimax rate is $q$; at $v_->B_0^2$ the
compatible-law field is empty.
\end{theorem}

\begin{proof}
The upper bound follows from
Proposition~\ref{prop:affine-histogram-adaptation}. For the lower, write
$v_0=v_-$ and use $m_\rho,L$ from (C.5). Put
\[
 \begin{gathered}
 \epsilon_\rho=\min\{1,c_{\rm rad},B_{\mathcal G}\}>0,\qquad
 p=2^{N+6},\qquad d=N+7,\\
 H_\mu=\epsilon_\rho s^2/128,\qquad
 H_\pi=\epsilon_\rho t^2/128.
 \end{gathered}
\]
Embed $(\mathbb Z/2\mathbb Z)^d$ measurably into the fixed infinite
standard Borel space, as in Appendix~\ref{app:p-raw}; denote its image
by $S$ and let $P_X^*$ be uniform on $S$. There are $2^d-1\ge p$
nonzero Walsh characters, so choose $p$ distinct characters
$\phi_1,\ldots,\phi_p$, extended by zero off $S$.

Publish the zero-centered full interval boxes with width $H_\nu$ on
every singleton $\{x\}$, $x\in S$, and width zero on $\mathcal X\setminus S$.
Rational normalized coordinates give a countable pointwise-everywhere
dense representation. This fixes $\Gamma^*=\Gamma^*(\rho,n,s,t)$
before any approximation budget, hard-family component, label,
coordinate, or sign is selected.

\paragraph{Both public complexity certificates.}
For either side with stochastic budget $r\in\{s,t\}$, the full
difference class has pointwise envelope $2H_\nu$. Its localized outer
Rademacher complexity is therefore at most $2H_\nu$ at every radius.
When $r>0$ and $R\ge r$,
\[
 \frac{\mathfrak R^*_{n,P_X^*}(R;\Delta\mathcal G_\nu)}{3R}
 \le\frac{2H_\nu}{3r}
 =\frac{\epsilon_\rho r}{192}\le c_{\rm rad}r.
\]
When $r=0$, that side is the singleton zero class, and its complexity
vanishes at every positive radius. Also
$H_\nu\le\epsilon_\rho/128\le B_{\mathcal G}$. These are certificates
for both full boxes in every family below, regardless of which
nuisance truth is represented.

\paragraph{Four legal families inside the same pair.}
Write $z=\sigma\phi_j$. We use the exact treatment channel (C.6),
which has the specified conditional mean and residual variance exactly
$v_0$ by (C.7). Given $(X,T)$, use the binary response channel (C.12)
with response amplitude $B_0/2$.

Two label templates suffice. The P template has parameters
$\eta,\gamma\ge0$ and
\[
 \begin{gathered}
 \Delta=\frac{2\eta\gamma}{v_0},\qquad
 \kappa=\gamma+\Delta\eta,\\
 (\beta,\pi,\mu)_0=(0,z\eta,z\gamma),\quad
 (\beta,\pi,\mu)_1=(\Delta,-z\eta,z\kappa).
 \end{gathered}
\]
The O template has parameters $\eta,\gamma\ge0$ and
\[
 \begin{gathered}
 \Delta=\frac{\gamma\eta}{v_0+\gamma^2},\\
 (\beta,\pi,\mu)_0=(0,z\gamma,z\eta),\quad
 (\beta,\pi,\mu)_1=(\Delta,z\gamma,0).
 \end{gathered}
\]
Their complete conditional masses are respectively (C.13) and (D.4).
Choose their amplitudes as follows:
\begin{equation}
\begin{array}{c|c|c|c}
 \text{term}&\text{template}&\eta&\gamma\\ \hline
 ab&{\rm P}&m_\rho b/4&a/512\\
 bs^2&{\rm P}&m_\rho b/4&\epsilon_\rho s^2/256\\
 at^2&{\rm O}&a/256&\epsilon_\rho m_\rho t^2/4\\
 s^2t^2&{\rm O}&\epsilon_\rho s^2/256&
                        \epsilon_\rho m_\rho t^2/4
\end{array}
\label{eq:adaptation-box-witness-scales}
\end{equation}
In both P rows,
$2\eta^2/v_0\le m_\rho^2/(8v_0)<1$, so $\kappa\le2\gamma$.
In the $ab$ row, the zero candidates give approximation errors at
most $\kappa\le a/256\le a$ and $\eta\le b$.
In the $bs^2$ row, both outcome truths belong to the outcome box
because $\kappa\le H_\mu$, while zero gives propensity error at most
$b$. In both O rows,
$\gamma\le\epsilon_\rho t^2/256\le H_\pi$, so propensity is represented.
The $at^2$ row has outcome error at most $\eta\le a$ using zero;
in the last row both outcome truths are represented since
$\eta\le H_\mu$. Exact approximation distances are unnecessary.

Every treatment mean has absolute value at most $m_\rho/4$.
The channel inequalities (C.5)--(C.7) ensure strict treatment
probabilities and $|T|,|u|<B_0$. Uniformly over these four rows,
\[
 0\le\Delta\le\frac1{2048B_0},\qquad |\mu|\le\frac1{128},
 \qquad |\beta T+\mu|\le\frac{17}{2048}<\frac1{64}.
\]
The binary response probabilities are consequently strict, and
$|\varepsilon|\le B_0/2+1/64<B_0$.
The remaining target, nuisance, and response envelopes follow from
these same bounds. Thus every component is a legal PLM, with exact
residual variance $v_0$, not merely a moment-matched array.

\paragraph{Common reference and mixture distance for every sample size.}
For the P template, the sign averages of the two labels coincide
atomwise by (C.15)--(C.16): their required identity is
\[
 \Delta(v_0+\eta^2)=\eta(\kappa+\gamma),
\]
which follows from $\kappa=\gamma+\Delta\eta$ and
$\Delta v_0=2\eta\gamma$. For the O template, (D.8)--(D.9) give the
same conclusion from $\Delta(v_0+\gamma^2)=\gamma\eta$.
Let $Q$ denote the resulting common law of the treatment and binary
response sign. It is strictly positive because it averages strictly
positive component masses. Only the component laws are required to
have exact residual variance $v_0$.

For either label, its two sign laws can be written
\[
 P_{\ell,z}=Q(1+z h_\ell),\qquad
 h_\ell=\frac{P_{\ell,+}-P_{\ell,-}}{2Q},\qquad
 \E_Qh_\ell=0,\quad |h_\ell|<1.
\]
Put $K_\ell=\E_Qh_\ell^2\le1$ and
$\mathbb Q=P_X^*\otimes Q$. A full one-observation likelihood is
$1+\sigma\phi_j(X)h_\ell(T,Y)$, where the binary sign is identified
with $Y/(B_0/2)$. Walsh orthogonality gives Gram entries
$1+K_\ell$, $1-K_\ell$, and one, as in Lemma~\ref{lem:testing-tools}.

Choose $(j,\sigma)$ once uniformly from $[p]\times\{-1,1\}$ and then
sample $N$ iid observations from that component. For each label's
resulting mixture, (C.3) gives
\[
 \chi^2(\overline P_\ell^{\,N},\mathbb Q^N)
 =\frac{\{(1+K_\ell)^N+(1-K_\ell)^N\}/2-1}{p}
 \le\frac{2^{N-1}-1}{2^{N+6}}<\frac1{128}.
\]
The cross-label total variation is at most $1/\sqrt{128}<1/4$.
Equation (C.1) therefore yields a lower risk of at least $\Delta/8$
for every row of \eqref{eq:adaptation-box-witness-scales}, for all
$n\ge1$. A zero term has zero
separation and causes no exception. The public pair is already fixed,
and all public evaluations are deterministic functions of $X$, so
disclosing the entire represented pair does not change this
full-observation comparison.

The four gaps are bounded below by their corresponding terms times
\[
 k_\rho=\frac{\epsilon_\rho^2m_\rho}
                {1024(v_0+m_\rho^2/16)}>0.
\]
Indeed the P gaps are
$m_\rho ab/(1024v_0)$ and
$\epsilon_\rho m_\rho bs^2/(512v_0)$; the O gaps are
$\epsilon_\rho m_\rho at^2/\{1024(v_0+\gamma^2)\}$ and
$\epsilon_\rho^2m_\rho s^2t^2/\{1024(v_0+\gamma^2)\}$.

\paragraph{Parametric term and assembly.}
Keep the same $P_X^*$ and make it independent of the scalar experiment
(C.8), with both nuisances zero, $T=\pm\sqrt{v_0}$ equiprobably, and
$\beta=\pm d_\rho^\star/\sqrt n$. The zero truths belong to both
published boxes, whose complexity certificates have already been
verified. The independent covariate factor leaves the calculation
leading to (C.9) unchanged, so this same fixed pair has lower risk
$3d_\rho^\star q/8$.

For every rule on $\Gamma^*$, its compatible-law supremum is therefore
at least
\[
 \frac18\min\{3d_\rho^\star,k_\rho\}
       \max\{q,ab,bs^2,at^2,s^2t^2\}.
\]
Taking the learner infimum and using that the maximum of five
nonnegative numbers is at least their sum divided by five proves the
lower bound with
\[
 c_{{\rm box},\rho}
 =\frac1{40}\min\{3d_\rho^\star,k_\rho\}>0.
\]
The construction fixes the same $\Gamma^*$ and $P_X^*$ for every
$(a,b)$; only the compatible hard laws vary. The outer supremum in
\eqref{eq:adaptation-box-minimax-risk} then gives the stated result.

At $v_-=B_0^2$, (H.3)--(H.4) give the $q$ upper bound for every
nonempty pair, and the zero singleton pair used in (H.5)--(H.6) is a
full interval box and supplies the $q$ lower bound. When $v_->B_0^2$,
the law field is empty by Appendix~\ref{app:boundary}; no numerical
risk is assigned to that empty field.
\end{proof}
